\documentclass[11pt,letterpaper]{article}
\usepackage[T1]{fontenc}
\usepackage[utf8]{inputenc}
\usepackage{newtxtext}
\usepackage[margin=1in,headheight=15pt]{geometry}
\usepackage{amsmath,amsthm,mathtools}
\usepackage{newtxmath}
\usepackage{microtype}
\usepackage{xcolor,booktabs,tabularx,array}
\usepackage{enumitem,fancyhdr,needspace}
\usepackage[most]{tcolorbox}
\usepackage{xurl}
\usepackage{hyperref}
\usepackage{aliascnt}
\usepackage[nameinlink,noabbrev]{cleveref}
\definecolor{Olive}{HTML}{4B5940}
\definecolor{Pale}{HTML}{F1F3EA}
\definecolor{Ink}{HTML}{253026}
\hypersetup{colorlinks=true,linkcolor=Olive,citecolor=Olive,urlcolor=Olive,
 pdftitle={Dilation Rigidity of Surreal Numbers and Omnific Integers},
 pdfauthor={Research article prepared with ChatGPT for Vladimir Reshetnikov}}
\pagestyle{fancy}\fancyhf{}
\fancyhead[L]{\small Dilation rigidity}
\fancyhead[R]{\small Surreal, surcomplex, and omnific numbers}
\fancyfoot[C]{\thepage}
\setlength{\parskip}{3pt}
\setlength{\emergencystretch}{2em}
\widowpenalty=10000\clubpenalty=10000
\allowdisplaybreaks[2]
\setcounter{tocdepth}{1}
\numberwithin{equation}{section}
\newtheorem{theorem}{Theorem}[section]
\newaliascnt{lemma}{theorem}
\newtheorem{lemma}[lemma]{Lemma}
\aliascntresetthe{lemma}
\crefname{lemma}{lemma}{lemmas}
\Crefname{lemma}{Lemma}{Lemmas}
\newaliascnt{proposition}{theorem}
\newtheorem{proposition}[proposition]{Proposition}
\aliascntresetthe{proposition}
\crefname{proposition}{proposition}{propositions}
\Crefname{proposition}{Proposition}{Propositions}
\newaliascnt{corollary}{theorem}
\newtheorem{corollary}[corollary]{Corollary}
\aliascntresetthe{corollary}
\crefname{corollary}{corollary}{corollaries}
\Crefname{corollary}{Corollary}{Corollaries}
\theoremstyle{definition}
\newaliascnt{definition}{theorem}
\newtheorem{definition}[definition]{Definition}
\aliascntresetthe{definition}
\crefname{definition}{definition}{definitions}
\Crefname{definition}{Definition}{Definitions}
\newaliascnt{example}{theorem}
\newtheorem{example}[example]{Example}
\aliascntresetthe{example}
\crefname{example}{example}{examples}
\Crefname{example}{Example}{Examples}
\newtheorem{question}{Research question}
\theoremstyle{remark}
\newaliascnt{remark}{theorem}
\newtheorem{remark}[remark]{Remark}
\aliascntresetthe{remark}
\crefname{remark}{remark}{remarks}
\Crefname{remark}{Remark}{Remarks}
\newcommand{\No}{\mathbf{No}}
\newcommand{\Oz}{\mathbf{Oz}}
\newcommand{\SC}{\mathbf{No}[i]}
\newcommand{\Goz}{\mathbf{Oz}[i]}
\newcommand{\kk}{k}
\newcommand{\C}{\mathbb C}
\newcommand{\R}{\mathbb R}
\newcommand{\Q}{\mathbb Q}
\newcommand{\Z}{\mathbb Z}
\newcommand{\N}{\mathbb N}
\newcommand{\A}{A_\Gamma(k)}
\newcommand{\K}{K_\Gamma(k)}
\newcommand{\val}{v}
\newcommand{\ct}{\operatorname{ct}}
\newcommand{\supp}{\operatorname{supp}}
\newcommand{\ord}{\operatorname{ord}}
\newcommand{\lc}{\operatorname{lc}}
\newcommand{\Res}{\operatorname{Res}}
\newcommand{\Frac}{\operatorname{Frac}}
\newcommand{\trdeg}{\operatorname{trdeg}}
\newcommand{\id}{\operatorname{id}}
\newcommand{\abs}[1]{\lvert#1\rvert}
\newcommand{\repo}{\texttt{VladimirReshetnikov/Surreal}}
\newcolumntype{Y}{>{\raggedright\arraybackslash}X}
\tcbset{colback=Pale,colframe=Olive,boxrule=.5pt,arc=1pt,
 left=9pt,right=9pt,top=8pt,bottom=8pt}

\begin{document}
\hypersetup{pageanchor=false}
\begin{titlepage}
\vspace*{0.15in}
{\large\scshape\color{Olive} Research article\par}
\vspace{0.22in}
{\Huge\bfseries Dilation Rigidity\par}
\vspace{0.10in}
{\LARGE\bfseries of Surreal Numbers\par}
\vspace{0.07in}
{\LARGE\bfseries and Omnific Integers\par}
\vspace{0.22in}
{\large Autonomous Mahler equations, finite support,\\
sharp algebraic degrees, and a higher-order hierarchy\par}
\vspace{0.25in}
{Research article prepared with ChatGPT\\for Vladimir Reshetnikov\par}
\vspace{0.08in}
{23 September 2026\par}
\vspace{0.25in}
\begin{abstract}
For an integer $d\ge2$, exponent dilation sends
$\sum a_\gamma t^\gamma$ to $\sum a_\gamma t^{d\gamma}$.
We prove that a nonconstant Hahn series algebraically dependent with its
own dilate over the constant field must belong to a single ordinary Laurent
series subfield. This holds in characteristic zero over every ordered
abelian value group, without an Archimedean or divisibility assumption.
If $F(y,S_dy)=0$, the necessary ramification is at most $\deg_YF$.
The proof combines Newton--Puiseux branches with a formal B\"ottcher
coordinate and a least-support-term argument; every evaluation is justified
by strong Hahn summability, not by a contraction principle.

For the negative-support ring this becomes an exact classification:
first-order autonomous algebraic dependence occurs precisely for polynomials
in one monomial. The least degree in the dilated variable is exactly the
degree in the primitive monomial. Consequently every infinite-support
omnific integer is algebraically independent of each of its nontrivial
integer dilates. We also obtain a finite-profile description for a fixed
algebraic equation, an exact Newton-weight existence test, complete
rational-map solutions, and a sharp infinitesimal-tail obstruction to
omnific solutions. Higher order behaves differently: an infinite-support
omnific integer has exact autonomous order two, and examples realize every
finite order. A research agenda identifies the remaining classification,
effectivity, positive-characteristic, and formalization problems.
\end{abstract}
\vfill
\begin{tcolorbox}
\small\textbf{Research status.} The proposed contributions are the
arbitrary-rank classifications and their consequences, not the classical
Newton--Puiseux or B\"ottcher theorems. Complete written arguments are supplied
relative to the stated classical inputs. The results have not been refereed
or Lean-verified, and priority is not certified. No named published conjecture
is claimed solved. Exact finite checks accompany the article but do not
verify the infinite theorems.
\end{tcolorbox}
\end{titlepage}
\hypersetup{pageanchor=true}
\pagenumbering{roman}
\tableofcontents
\clearpage
\pagenumbering{arabic}

\section{Research position and the questions answered}\label{sec:position}

An omnific integer can have a transfinite normal form, infinitely many
nonzero coefficients, and exponents at mutually incomparable Archimedean
scales. A polynomial in one monomial has none of this support complexity.
The principal question of this article is whether an algebraic relation
between a number and its own exponent-dilated copy can distinguish these
situations.

The dilation is
\begin{equation}\label{eq:dilationintro}
 S_d\!\left(\sum_\alpha a_\alpha\omega^{g_\alpha}\right)
   =\sum_\alpha a_\alpha\omega^{d g_\alpha}.
\end{equation}
It is a field automorphism on $\No$ and $\SC$, and fixes ordinary constants.
It is \emph{not} the map $y\mapsto y^d$. For example,
\[
 S_2(\omega+1)=\omega^2+1,
 \qquad (\omega+1)^2=\omega^2+2\omega+1.
\]
The question is about polynomial equations over the ordinary constant
field, with $S_d$ additionally available. It is not a statement that an
infinite omnific integer is algebraic over $\R$.

\subsection{The main answer in omnific notation}

For $x\in\Oz\setminus\Z$ and $d\ge2$, we prove the equivalence
\begin{equation}\label{eq:headline}
 \begin{split}
 &x\text{ and }S_dx\text{ are algebraically dependent over }\R\\
 &\quad\Longleftrightarrow\quad
 x=P(\omega^\delta)\text{ for some }\delta>0,
 \quad P\in\R[T],\quad P(0)\in\Z.
 \end{split}
\end{equation}
The same statement holds for Gaussian omnific integers with $\C$ and
$\Z[i]$ in place of $\R$ and $\Z$. Thus an infinite normal-form support is
an unconditional obstruction to \emph{every} first-order autonomous
algebraic relation. A finite support can also obstruct a relation: its
positive exponents must be commensurable, not merely comparable in size.

There is an exact quantitative refinement. Normalize $\delta$ so that the
nonzero positive exponents of $x$ generate $\delta\Z$. Then, writing
$x=P(\omega^\delta)$ and $M=\deg P$,
\begin{equation}\label{eq:headline-degree}
 \min\{\deg_YF:0\ne F\in\R[X,Y],\ F(x,S_dx)=0\}=M.
\end{equation}
In particular, the answer is independent of the chosen integer $d\ge2$.
This is a classification of a substantial class of equations, rather than
one specially selected identity.

\subsection{Relation to the repository}

The repository was inspected at commit
\begin{center}\small\texttt{b895e8672990e8b5a56f97dd7246f9dd5f86c771}.\end{center}
The main points of contact are the report on a single dilation
\cite{RepoDilation}, the dynamics report \cite{RepoDynamics}, and the
omnific and Hahn-support infrastructure catalogued in the repository
\cite{Repo}. In particular, the single-dilation report already uses the
recognizer $S_2w=w^2$ for canonical monomials and treats linear and
multiplicative difference equations. That recognizer is a premise to
extend, not a novelty to reclaim. We reprove its elementary
characteristic-zero version for general $d$.

The present extension is from this one special equation to
\emph{every} constant-coefficient polynomial relation $F(y,S_dy)=0$,
including ramified algebraic correspondences. It supplies support bounds,
an exact omnific classification, and sharp degrees. The inspected
repository excerpts do not supply this classification. That comparison
is not an assertion that every line of every repository manuscript has
been independently reviewed.

\subsection{Literature and the novelty boundary}

Hahn summation and its support calculus are classical \cite{Neumann}.
Newton--Puiseux factorization is classical; Mannaa--Coquand give a modern
constructive treatment \cite{MC}. Formal B\"ottcher coordinates are also
classical; the needed characteristic-zero theorem is recorded explicitly
by Salerno--Silverman \cite{SS}. We use these inputs openly and give an
explicit coordinate construction adapted to strong Hahn evaluation.

There is an established literature on autonomous Mahler equations,
including Nishioka--Nishioka \cite{NN}, and on nonlinear formal Mahler
solutions \cite{GG}. Chyzak--Dreyfus--Dumas--Mezzarobba \cite{CDDM} and
Faverjon--Roques \cite{FR} address linear Mahler equations and generalized
series. The later works of Faverjon--Poulet \cite{FP} and
Faverjon--Roques \cite{FRPurity} develop basis algorithms and generalized
Mahler-series structure for linear systems; they are cited here for
context from their abstracts, not used as proof dependencies.
These works prevent any claim that Mahler equations, B\"ottcher
coordinates, or non-Puiseux solutions are new topics. The contribution
proposed here is the \emph{arbitrary-rank support-collapse theorem}, the
resulting exact omnific criterion, and their quantitative synthesis.

For omnific arithmetic, L'Innocente--Mantova \cite{LM} provide the relevant
published factorization framework. Our result is not a factorization theorem
and does not settle the remaining general factorization conjectures.
The user's encyclopedia link \cite{Wiki} is useful orientation, but none
of the proofs relies on an encyclopedia assertion in place of the
classical inputs above.

\section{Hahn fields, normal forms, and the two support directions}
\label{sec:setup}

\begin{definition}
Let $k$ be a field and $\Gamma$ an ordered abelian group. The full Hahn
field $\K=k((t^\Gamma))$ consists of formal sums
\[
 y=\sum_{\gamma\in\Gamma}a_\gamma t^\gamma,
 \qquad a_\gamma\in k,
\]
whose support is a well-ordered subset of $\Gamma$. For $y\ne0$, set
$v(y)=\min\supp y$ and let $\lc(y)$ be the coefficient there. Put
$v(0)=+\infty$.
\end{definition}

Throughout the main classification, $\operatorname{char}k=0$ and $d$ is
an ordinary integer at least two. No order on $k$ is required. Dilation
\begin{equation}\label{eq:Sd}
 S_d\Bigl(\sum a_\gamma t^\gamma\Bigr)=
 \sum a_\gamma t^{d\gamma}
\end{equation}
is a strong injective field endomorphism; it is an automorphism when
$\Gamma$ is $d$-divisible. Its elementary valuation law is
\begin{equation}\label{eq:val-dilation}
 v(S_dy)=d\,v(y).
\end{equation}

The valuation ring and its maximal ideal are
\[
 \mathcal O=\{y:v(y)\ge0\},\qquad
 \mathfrak m=\{y:v(y)>0\}.
\]
For $y\in\mathcal O$, its residue is its coefficient at zero, an element
of $k$. In contrast, define
\begin{equation}\label{eq:A}
 \A=\{y\in\K:\supp y\subseteq\Gamma_{\le0}\}.
\end{equation}
This is a subring, but it is \emph{not} the valuation ring. A nonconstant
element of $\A$ has negative valuation. The absence of \emph{positive}
exponents, not just a bound on the leading exponent, will be decisive.

\begin{definition}
For $\delta\in\Gamma_{>0}$, the notation $k((t^\delta))$ means the image
of the ordinary Laurent series field $k((Z))$ under $Z\mapsto t^\delta$.
Thus its elements have support contained in $\delta\Z$, bounded below
in the integer indexing. A series with this property is called
\emph{cyclically supported} here.
\end{definition}

Cyclic support is stronger than support in one Archimedean class and
stronger than a one-dimensional rational span. The set
$\{-1,-1/2,-1/4,\ldots\}$ spans a one-dimensional $\Q$-space but is not
contained in any cyclic subgroup of $\Q$.

\subsection{Translation to surreal numbers}

Conway normal forms use decreasing powers of $\omega$. The correspondence
with the increasing Hahn convention is
\begin{equation}\label{eq:omega-bridge}
 t^\gamma\longmapsto\omega^{-\gamma}.
\end{equation}
For any set-sized ordered additive subgroup $\Gamma\subseteq\No$, it
identifies $\R((t^\Gamma))$ with the corresponding normal-form subfield
of $\No$; complexifying gives $\C((t^\Gamma))\subseteq\SC$.
Every individual surreal or surcomplex number is contained in such a
workspace: take the divisible subgroup generated by its set-sized support.
The usual normal-form arithmetic is the algebraic input here
\cite{LM,Repo}.

The omnific integers have the description
\begin{equation}\label{eq:Oz}
 \Oz=\left\{n+\sum_{\alpha<\lambda}r_\alpha\omega^{g_\alpha}:
 n\in\Z,\ r_\alpha\in\R,\ g_\alpha>0,
 \ (g_\alpha)_{\alpha<\lambda}\text{ strictly decreasing}\right\},
\end{equation}
with set-sized support and zero coefficients omitted. Thus the only
additional restriction beyond the real negative-support ring is the
integrality of the coefficient at zero. For $\Goz$ the coefficients are
complex and that constant lies in $\Z[i]$.

\subsection{Size and topology conventions}

All field-theoretic proofs below take place in a set-sized Hahn field.
The statements for $\No$ and $\SC$ are then elementwise consequences.
No proper-class polynomial ring, class-sized tensor product, or
class-indexed analytic sum is required.

Also, $v(x)>0$ need not imply that $v(x^n)$ tends cofinally to infinity
in $\Gamma$. For example, in the lexicographically ordered group
$\Q\oplus\Q$, the multiples of $(0,1)$ never reach $(1,0)$.
Nevertheless the geometric series in $x$ is strongly summable. This
separation between support summability and topological convergence is
essential throughout.

\section{The support calculus and a monomial recognizer}
\label{sec:support}

A family $(f_i)$ of Hahn series is \emph{strongly summable} when the union
of its supports is well ordered and every exponent occurs in the support
of only finitely many family members. Its sum is then defined coefficientwise.

\begin{lemma}[Formal evaluation]\label{lem:evaluation}
Let $x\in\mathfrak m$. Every formal series $f(Z)\in k[[Z]]$ can be evaluated
at $x$ as a strong Hahn sum. This gives a $k$-algebra homomorphism.
For $x\ne0$ it extends to an injective map $k((Z))\to\K$, and
\begin{equation}\label{eq:evaluation-value}
 v(f(x))=\ord_Z(f)\,v(x)\qquad(f\ne0).
\end{equation}
Formal composition is respected whenever the inner series has zero
constant coefficient. Moreover,
\begin{equation}\label{eq:commute-eval}
 S_d(f(x))=f(S_dx).
\end{equation}
\end{lemma}
\begin{proof}
The Neumann support lemma says that the monoid generated by a well-ordered
subset of $\Gamma_{>0}$ is well ordered and that each exponent has only
finitely many representations by finite words from that subset, allowing
all word lengths \cite{Neumann}. Apply it to $\supp x$. It justifies the
family of powers and the regroupings for products and formal composition.
There are only finitely many negative powers in a Laurent series, so these
cause no extra infinite-sum issue.

If the initial term of $f$ is $aZ^m$, its evaluation has initial valuation
$m\,v(x)$: every later power has strictly larger valuation. This proves
\eqref{eq:evaluation-value} and injectivity. Exponent dilation preserves
well ordering, finite occurrence, and coefficientwise sums, proving
\eqref{eq:commute-eval}.
\end{proof}

\begin{lemma}[Constants are relatively algebraically closed]
\label{lem:constants}
Every element of $\K$ algebraic over $k$ belongs to $k$. In particular,
every nonconstant Hahn series is transcendental over $k$, and
$\operatorname{Fix}(S_d)=k$.
\end{lemma}
\begin{proof}
A nonzero algebraic element $z$ satisfies a polynomial with nonzero
constant and leading coefficients in $k$. If $v(z)>0$, the constant term
is the unique term of least valuation; if $v(z)<0$, the highest-degree
term is. Both are impossible. Hence $v(z)=0$. Let $c\in k$ be its
residue. If $z-c\ne0$, this is an algebraic element of positive valuation,
again impossible.

For a nonconstant fixed element $y$, use $x=y-c$ if $v(y)\ge0$, with
$c$ its residue, and $x=1/y$ if $v(y)<0$. In both cases $x\ne0$ has
positive valuation and is fixed. The equation $d\,v(x)=v(x)$ contradicts
$d>1$ in an ordered group.
\end{proof}

\begin{lemma}[Monomial recognizer]\label{lem:monomial}
Suppose $w\ne0$ and $S_dw=w^d$. Then
\begin{equation}\label{eq:monomial-recognizer}
 w=a t^\eta,\qquad a\in k^\times,\quad a^{d-1}=1,
 \quad\eta\in\Gamma.
\end{equation}
Conversely every such monomial is a solution.
\end{lemma}
\begin{proof}
Write $w=a t^\eta(1+\varepsilon)$ with $\varepsilon\in\mathfrak m$.
Initial coefficients give $a=a^d$. Dividing the original identity by
$a t^{d\eta}$ gives
\[
 1+S_d\varepsilon=(1+\varepsilon)^d.
\]
If $\varepsilon\ne0$ and $\rho=v(\varepsilon)>0$, then the nonconstant
part of the left side has valuation $d\rho$, whereas that of the right
side has valuation $\rho$: its initial coefficient is multiplied by the
nonzero scalar $d$. Since $d\rho>\rho$, this is impossible.
\end{proof}

The least-term proof eliminates possible tails at \emph{all} other scales
at once. It neither examines the terms sequentially nor assumes that an
iteration approaches its limit in the ambient valuation topology.
The $d=2$ instance is the repository's monomial recognizer
\cite{RepoDilation}.

\section{The classical ramification input and its degree bound}
\label{sec:puiseux}

We use the Newton--Puiseux theorem in the following standard form:
if $k$ is algebraically closed of characteristic zero, then
\[
 \bigcup_{e\ge1}k((U^{1/e}))
\]
is an algebraic closure of $k((U))$ \cite{MC}. Thus every nonzero
polynomial in $k((U))[V]$ splits into linear factors there. The next
lemma spells out both the quantitative point and its transport into a
Hahn field.

\begin{lemma}[Evaluated Puiseux branch]\label{lem:puiseux}
Assume $k$ algebraically closed, $\Gamma$ divisible, and $x\in\K$ with
$v(x)=\gamma>0$. Let $0\ne G(U,V)\in k[U,V]$ have $V$-degree $D$, and
suppose $G(x,z)=0$. Then there exist an integer $1\le e\le D$, an
$e$th root $u$ of $x$ in $\K$, and $\phi(Z)\in k((Z))$ such that
\begin{equation}\label{eq:puiseux-branch}
 z=\phi(u),\qquad G(Z^e,\phi(Z))=0.
\end{equation}
Moreover $v(u)=\gamma/e$, and
$v(z)=\ord_Z(\phi)\,\gamma/e$ when $z\ne0$.
\end{lemma}
\begin{proof}
Apply Newton--Puiseux to $G(U,V)$ over $k((U))$. In a common extension
$k((U^{1/N}))$ it factors as a nonzero leading coefficient times finitely
many factors $V-\psi_j(U^{1/N})$.

An $N$th root of $x$ exists in the Hahn field without using its full
algebraic closedness: write $x=a t^\gamma(1+\varepsilon)$, choose
$b^N=a$, and use
\[
 b t^{\gamma/N}(1+\varepsilon)^{1/N},
\]
where the binomial series is evaluated by \cref{lem:evaluation}.
Evaluation at this root sends the Puiseux factorization to a genuine
factorization in $\K$. Its leading coefficient is nonzero by the
injectivity in \cref{lem:evaluation}. Consequently $z$ equals one of
the evaluated roots.

Choose the smallest denominator $e$ for that root. Its series can be
written as $\phi(U^{1/e})$, with $e\mid N$; the corresponding power of
the chosen $N$th root gives $u$. Minimality means that the automorphisms
$U^{1/e}\mapsto\zeta U^{1/e}$, for $\zeta^e=1$, give exactly $e$
distinct conjugates of this series. Indeed a stabilizing $\zeta$ must
satisfy $\zeta^j=1$ for every exponent $j$ occurring in $\phi$, and
minimality says their common divisor with $e$ is one. All these conjugates
are roots of $G$. Hence $e\le D$. The valuation statements follow from
\cref{lem:evaluation}.
\end{proof}

\begin{remark}
The bound is on the denominator of \emph{one selected branch}, not on a
common denominator for every root. A common splitting denominator can be
larger than $D$. Confusing these two denominators would invalidate the
sharp degree theorem below.
\end{remark}

\section{A support-safe B\"ottcher construction}\label{sec:bottcher}

The formal-coordinate theorem in this section is classical
\cite[Proposition~1]{SS}. We include a construction so that no unstated
convergence hypothesis enters its Hahn-field application. All sums in
the construction are first formed in the ordinary formal ring $k[[Z]]$.

\begin{lemma}[Normalized formal coordinate]\label{lem:bottcher}
Let $h(Z)=Z^dU(Z)$, where $U\in1+Zk[[Z]]$ and
$\operatorname{char}k=0$. There is a unique
$b(Z)\in Z+Z^2k[[Z]]$ with
\begin{equation}\label{eq:bottcher-b}
 b(h(Z))=b(Z)^d.
\end{equation}
Its compositional inverse $H$ satisfies
\begin{equation}\label{eq:bottcher-H}
 H(Z^d)=h(H(Z)).
\end{equation}
\end{lemma}
\begin{proof}
Put
\begin{equation}\label{eq:bottcher-log}
 L(Z)=\sum_{n\ge0}\frac{1}{d^{n+1}}
       \log U\bigl(h^{\circ n}(Z)\bigr),\qquad
 b(Z)=Z\exp L(Z).
\end{equation}
The $n$th summand has order at least $d^n$ unless it is zero, so each
coefficient receives only finitely many contributions. Reindexing gives
\[
 dL(Z)=\log U(Z)+L(h(Z)).
\]
Therefore
\[
 b(h(Z))=Z^dU(Z)\exp L(h(Z))=Z^d\exp(dL(Z))=b(Z)^d.
\]
The linear coefficient of $b$ is one, so it has a unique compositional
inverse, which gives \eqref{eq:bottcher-H}.

For uniqueness, any normalized coordinate has the form $Z\exp M(Z)$.
The difference $D=L-M$ satisfies $D(h)=dD$. A nonzero $D$ of positive
order would have different orders on these two sides, namely
$d\ord D$ and $\ord D$. Thus $D=0$.
\end{proof}

\begin{proposition}[All Hahn solutions of a superattracting germ]
\label{prop:germ}
Let $h(Z)=A Z^d+O(Z^{d+1})\in k[[Z]]$, $A\ne0$, and let $x\ne0$
have positive valuation. Then $S_dx=h(x)$ is solvable with prescribed
leading coefficient $a\in k^\times$ and valuation $\delta\in\Gamma_{>0}$
if and only if
\begin{equation}\label{eq:leading-a}
 A a^{d-1}=1.
\end{equation}
For each such pair there is exactly one solution:
\begin{equation}\label{eq:germ-profile}
 x=aH_a(t^\delta),
\end{equation}
where $H_a$ is the normalized inverse B\"ottcher coordinate of
$h_a(Z)=a^{-1}h(aZ)$.
\end{proposition}
\begin{proof}
Leading coefficients force \eqref{eq:leading-a}. With this condition,
$h_a(Z)=Z^d+O(Z^{d+1})$. Write $u=x/a$, and let $b_a$ be the coordinate
of \cref{lem:bottcher}. Strong evaluation and its commutation with $S_d$
give
\[
 S_d(b_a(u))=b_a(S_du)=b_a(h_a(u))=b_a(u)^d.
\]
The left-hand argument $b_a(u)$ has valuation $\delta$ and leading
coefficient one. By \cref{lem:monomial} it is exactly $t^\delta$.
Inversion yields \eqref{eq:germ-profile} and uniqueness. Conversely,
\eqref{eq:bottcher-H}, evaluated strongly at $t^\delta$, verifies that
this expression is a solution.
\end{proof}

This is an exhaustive statement in the full Hahn field. The coordinate
series is ordinary, but the element to which it is applied may initially
have arbitrary support. The recognizer proves that every extra support
scale disappears after the coordinate change.

\section{The arbitrary-rank support-collapse theorem}\label{sec:main}

\begin{theorem}[First-order autonomous support rigidity]
\label{thm:collapse}
Let $k$ be any field of characteristic zero, $\Gamma$ any ordered abelian
group, and $d\ge2$. Suppose $y\in\K\setminus k$ and
\begin{equation}\label{eq:F-main}
 F(y,S_dy)=0,\qquad 0\ne F\in k[X,Y].
\end{equation}
Define a positive-valuation coordinate
\begin{equation}\label{eq:coordinate}
 x=\begin{cases}
 y-c,&v(y)\ge0,\quad c=\operatorname{res}(y),\\
 1/y,&v(y)<0,
 \end{cases}
 \qquad \gamma=v(x)>0.
\end{equation}
Then for some integer $e$ with $1\le e\le\deg_YF$, one has
\begin{equation}\label{eq:collapse}
 \delta=\gamma/e\in\Gamma,
 \qquad y\in k((t^\delta)).
\end{equation}
In particular, the support is finite or has order type $\omega$ and is
contained in a cyclic subgroup of $\Gamma$.
\end{theorem}

\begin{proof}
First assume $k$ algebraically closed and $\Gamma$ divisible. Let
$D=\deg_YF$. Since $y$ is transcendental over $k$, $D\ge1$.
The equation becomes $G(x,S_dx)=0$ for a nonzero polynomial $G$ with
$\deg_VG\le D$. In the finite case use
\[
 G(U,V)=F(c+U,c+V).
\]
In the infinite case, if $a=\deg_XF$, use
\[
 G(U,V)=U^aV^D F(U^{-1},V^{-1});
\]
monomial factors can be removed. These substitutions give nonzero
polynomials, and $x,S_dx$ are nonzero.

Apply \cref{lem:puiseux}. There are $1\le e\le D$, an element $z$ with
$z^e=x$, and a Laurent series $\phi$ such that
\begin{equation}\label{eq:mainbranch}
 S_dx=\phi(z),\qquad v(z)=\gamma/e.
\end{equation}
But $v(S_dx)=d\gamma$. Consequently
\[
 \ord_Z\phi=de,
 \qquad \phi(Z)=B Z^{de}(1+u(Z)),
 \quad B\ne0,\quad u\in Zk[[Z]].
\]
Choose an $e$th root of $B$ and form a formal $e$th root of this series.
Every such root has the form
\[
 h(Z)=A Z^d(1+u(Z))^{1/e}\in k[[Z]],\qquad A^e=B.
\]
Since $(S_dz)^e=S_dx=\phi(z)$, the ratio between $S_dz$ and one evaluated
root $h(z)$ is an $e$th root of unity. Multiplying $A$ by that root of
unity gives an $h$ with the exact identity
\begin{equation}\label{eq:z-germ}
 S_dz=h(z).
\end{equation}
By \cref{prop:germ},
$z=aH(t^{\gamma/e})$ for an ordinary formal power series $H$ and a
nonzero constant $a$. Thus $x=z^e$ belongs to
$k((t^{\gamma/e}))$, as does $y=c+x$ or $y=1/x$. This proves the
assertion in the auxiliary algebraically closed, divisible setting.

For general $k$ and $\Gamma$, embed the Hahn field coefficientwise and
exponentwise into
\[
 \overline{k}((t^{\Gamma_{\Q}})),
 \qquad \Gamma_{\Q}=\Gamma\otimes_{\Z}\Q
\]
with its natural ordered divisible hull. The proof gives an integer
$e\le D$ and support contained in $(\gamma/e)\Z$. The intersection
\[
 \Gamma\cap(\gamma/e)\Z=m(\gamma/e)\Z
\]
for a positive integer $m$ dividing $e$: it is a nonzero subgroup of a
cyclic group and contains $\gamma$. Put $e'=e/m$. Then
$1\le e'\le D$, $\gamma/e'\in\Gamma$, and the original support is
contained in $(\gamma/e')\Z$. Its coefficients are still in $k$, by
uniqueness of Hahn coefficients. The support is bounded below in that
cyclic group. Renaming $e'$ as $e$ proves \eqref{eq:collapse}.
\end{proof}

\begin{remark}[What the theorem does not assert]
Membership in a cyclic Laurent subfield is a necessary condition for
\eqref{eq:F-main}, not a claim that every ordinary Laurent series satisfies
some autonomous Mahler equation. Nor does the support conclusion turn
strong evaluation into ambient valuation convergence. The exponent group
can have arbitrarily high rank even when the individual solution is
forced into one cyclic subgroup.
\end{remark}

\section{Finite profiles, coefficient descent, and an existence test}
\label{sec:profiles}

The proof gives more than a necessary support condition. For a fixed
polynomial equation it supplies a finite collection of ordinary formal
shapes; the only remaining continuous scale parameter is the monomial
at which a shape is evaluated.

\begin{theorem}[Finite-profile classification]\label{thm:profiles}
Let $k$ be algebraically closed of characteristic zero, $d\ge2$, and
$0\ne F\in k[X,Y]$. There is a finite, possibly empty collection
\[
 \mathscr P(F,d)\subset k((Z))\setminus k
\]
such that for every nonzero divisible ordered group $\Gamma$,
\begin{equation}\label{eq:profiles}
 \{y\in\K\setminus k:F(y,S_dy)=0\}
  =\{f(t^\delta):f\in\mathscr P(F,d),\ \delta\in\Gamma_{>0}\}.
\end{equation}
Each profile is constructed from a Puiseux branch and a normalized
B\"ottcher coordinate. A profile at infinity has pole order at most
$\deg_YF$. The list need not be free of duplicates.
\end{theorem}
\begin{proof}
Remove all factors $Y-X$ from $F$, obtaining $F_0$. A nonconstant solution
cannot satisfy $S_dy=y$ by \cref{lem:constants}, so this removal does not
change the nonconstant solution set. The polynomial $F_0(X,X)$ is not
identically zero. Consequently there are only finitely many possible
finite residues $c$, namely its roots.

For each such $c$, form $G(U,V)=F_0(c+U,c+V)$. Also form the reciprocal
polynomial at infinity as in the proof of \cref{thm:collapse}. For each
of these finitely many polynomials, choose its finitely many Puiseux roots
$V=\phi(U^{1/e})$. Keep only the nonzero branches satisfying
$\ord_Z\phi=de$. For each kept branch, there are finitely many formal
$e$th roots
\[
 h(Z)^e=\phi(Z),\qquad h(Z)=A Z^d+O(Z^{d+1}),
\]
and finitely many constants $a$ with $Aa^{d-1}=1$. Let $H_a$ be the
coordinate in \cref{prop:germ}. Put on the list
\begin{equation}\label{eq:profile-shapes}
 f(Z)=c+(aH_a(Z))^e
 \quad\text{or}\quad
 f(Z)=(aH_a(Z))^{-e},
\end{equation}
according to the finite or infinite chart.

The proof of \cref{thm:collapse} shows that every nonconstant solution
has one of these forms after evaluation at a positive monomial.
Conversely let $z=aH_a(t^\delta)$. Then $S_dz=h(z)$, so with $x=z^e$,
\[
 S_dx=h(z)^e=\phi(z),\qquad G(x,S_dx)=0.
\]
Undoing the chart gives the original equation. The pole-order bound is
the branch denominator bound $e\le\deg_YF$.
\end{proof}

\begin{corollary}[Coefficient descent]\label{cor:coefficients}
Let $k_0\subseteq L$ be characteristic-zero fields and
$0\ne F\in k_0[X,Y]$. If $y\in L((t^\Gamma))\setminus L$ satisfies
$F(y,S_dy)=0$, then every Hahn coefficient of $y$ is algebraic over $k_0$.
In particular, an equation over $\overline{\Q}$ admits no nonconstant
solution having a coefficient transcendental over $\overline{\Q}$.
\end{corollary}
\begin{proof}
Work in an algebraic closure $\overline L$ and let $k'$ be the algebraic
closure of $k_0$ inside $\overline L$. Construct the finite profiles over
$k'$. Their Puiseux factorizations still split the relevant polynomials
over $\overline L((U))$; the roots of unity and leading-coefficient roots
are already in $k'$. Thus the same exhaustion proof gives profiles with
all coefficients in $k'$, even for a solution originally in the larger
coefficient field. Evaluation at a monomial changes exponents, not
coefficients. The original coefficients lie in $L\cap k'$.
\end{proof}

This is a descent statement for a \emph{fixed coefficient field of the
equation}. It does not contradict the existence of $a\omega^\delta$ as
a first-order solution over $\R$ for arbitrary real $a$: the corresponding
equation generally contains $a$ among its coefficients.

\begin{theorem}[Exact local Newton-weight criterion]\label{thm:newton}
Let $k$ be algebraically closed of characteristic zero and $\Gamma\ne0$
divisible. For
\[
 0\ne G(U,V)=\sum g_{ij}U^iV^j\in k[U,V],
\]
there is a nonzero $x\in\mathfrak m$ satisfying $G(x,S_dx)=0$ if and only if
\begin{equation}\label{eq:newton-weight}
 \min\{i+dj:g_{ij}\ne0\}
 \quad\text{is attained by at least two monomials.}
\end{equation}
\end{theorem}
\begin{proof}
If $v(x)=\gamma>0$, the valuation of a term is $(i+dj)\gamma$.
A unique least term cannot cancel, proving necessity.

For sufficiency use the ordinary Newton polygon, or equivalently the
Puiseux factorization. The weighted initial polynomial
\[
 I(W)=\sum_{i+dj=w}g_{ij}W^j,
 \qquad w=\min(i+dj),
\]
has at least two terms. Distinct active monomials have distinct $j$.
After its monomial factor is removed, it has a nonzero constant term and
positive degree, hence a nonzero root over $k$.
Factoring $G$ into Puiseux linear factors shows that these nonzero roots
are exactly the leading coefficients of the branches whose $U$-order is
$d$. More explicitly, a factor with root order below $d$ contributes a
nonzero constant to the weighted initial product, one above $d$ contributes
$W$, and one equal to $d$ contributes $W-A$ with $A\ne0$. The product
has a nonzero root precisely in the last case.

Thus some branch is $V=\phi(U^{1/e})$ with $\ord_Z\phi=de$.
The construction in \cref{thm:profiles} gives a nonzero positive-valuation
solution at every chosen positive monomial scale.
\end{proof}

\begin{corollary}[Scale-independent existence]\label{cor:existence}
For a fixed $F$ over an algebraically closed characteristic-zero field,
existence of a nonconstant solution is the same in every nonzero divisible
Hahn value group. It is decided by finitely many tests
\eqref{eq:newton-weight}: the finite diagonal-residue charts after removal
of $Y-X$ factors, and the chart at infinity. For algebraic-number
coefficients this is an effective decision procedure.
\end{corollary}
\begin{proof}
The charts and their number are as in \cref{thm:profiles}.
Each chart has a solution exactly when \cref{thm:newton} applies.
The test uses polynomial arithmetic, roots of one ordinary polynomial,
and comparison of finitely many integer weights. With exact algebraic
coefficients, all the required equality tests and root calculations are
effective. A nonzero divisible ordered group always has a positive element.
\end{proof}

The criterion decides \emph{existence}, not whether an arbitrary encoded
infinite Hahn series satisfies the equation. It also does not make an
arbitrary complex coefficient into computable input.

\section{The exact classification in the negative-support ring}
\label{sec:negative}

\begin{theorem}[Polynomial-in-one-monomial criterion]
\label{thm:negative}
Let $\operatorname{char}k=0$, $d\ge2$, and $y\in\A\setminus k$.
The following are equivalent:
\begin{enumerate}[label=\textup{(\roman*)},leftmargin=2em,itemsep=2pt]
\item $y$ and $S_dy$ are algebraically dependent over $k$;
\item $y=P(t^{-\delta})$ for some $\delta\in\Gamma_{>0}$ and $P\in k[T]$;
\item the support of $y$ is finite and its nonzero exponents generate a
      cyclic subgroup of $\Gamma$.
\end{enumerate}
If a relation $F(y,S_dy)=0$ is given, the polynomial in \textup{(ii)}
can be chosen of degree at most $\deg_YF$.
\end{theorem}
\begin{proof}
Since $y\notin k$ and its support has no positive exponent, $v(y)<0$.
Put $\gamma=-v(y)$. By \cref{thm:collapse}, for some $e\le\deg_YF$ its
support lies in $(\gamma/e)\Z$. It also lies between $-\gamma$ and zero.
There are only $e+1$ points of this cyclic subgroup in that interval.
Consequently
\[
 y=\sum_{j=0}^e a_jt^{-j\gamma/e},
\]
which proves (i)$\Rightarrow$(ii) with the degree bound.

For (ii)$\Rightarrow$(i), write $T=t^{-\delta}$. Then
$y=P(T)$ and $S_dy=P(T^d)$ belong to the one-variable rational function
field $k(T)$, whose transcendence degree is one. More explicitly, the
nonzero polynomial
\begin{equation}\label{eq:resultant}
 R_P(X,Y)=\Res_T(P(T)-X,\ P(T^d)-Y)
\end{equation}
vanishes on $(y,S_dy)$. Its degree in $Y$ is $\deg P$, so it is not
identically zero. Conditions (ii) and (iii) are equivalent by taking a
positive generator of the subgroup generated by the finitely many support
exponents.
\end{proof}

\begin{corollary}[Dilation independence]\label{cor:d-independent}
For $y\in\A$, dependence of $y,S_dy$ for one integer $d\ge2$ is
equivalent to dependence for every such integer. If $y$ has infinite
support, every pair $y,S_dy$ is algebraically independent over $k$.
The same independence conclusion holds for a finite support of rational
rank at least two.
\end{corollary}
\begin{proof}
The support condition in \cref{thm:negative} does not mention $d$.
For a finite subset of a torsion-free abelian group, the generated subgroup
is cyclic exactly when its rational span has dimension at most one.
\end{proof}

\begin{remark}[Why the support restriction is indispensable]
A Laurent series can have infinitely many positive exponents even when
it is cyclically supported. Intersecting with $\A$ cuts off all of them.
Thus the theorem is not a cardinality argument. Its force comes from a
\emph{discreteness conclusion} followed by a \emph{two-sided support bound}.
A merely bounded, nondiscrete support can still be infinite.
\end{remark}

\section{Sharp algebraic degrees and recovery of a monomial}
\label{sec:degree}

\begin{definition}
For a nonconstant finite-support $y\in\A$ satisfying the criterion, let
$\delta>0$ be the positive generator of the group generated by its
nonzero exponents. Write $y=P(t^{-\delta})$. The polynomial $P$ is
\emph{primitive in its exponents}: the greatest common divisor of its
positive supported degrees is one. Define $M(y)=\deg P$.
\end{definition}

\begin{theorem}[The exact relation degree]\label{thm:degree}
For every $d\ge2$,
\begin{equation}\label{eq:exact-degree}
 \min_{\substack{0\ne F\in k[X,Y]\\F(y,S_dy)=0}}\deg_YF=M(y).
\end{equation}
Moreover, if $T=t^{-\delta}$ is the primitive monomial, then
\begin{equation}\label{eq:recover}
 k(y,S_dy)=k(T).
\end{equation}
The irreducible relation has bidegree
$(\deg_X,\deg_Y)=(dM(y),M(y))$.
\end{theorem}
\begin{proof}
Let $M=M(y)$ and let a relation of $Y$-degree $D$ be given.
The degree-bounded conclusion of \cref{thm:negative} supplies
$y=Q(t^{-\delta'})$ of degree $e\le D$.
The support-generated subgroup $\delta\Z$ is contained in
$\delta'\Z$, so $\delta=m\delta'$ for an integer $m\ge1$.
Comparison of the largest support magnitude gives
$M\delta=e\delta'$, hence $e=Mm\ge M$. Therefore $D\ge M$.
The resultant \eqref{eq:resultant} gives the reverse inequality.

For the field statement, the extension
$k(T)/k(P(T))$ has degree $M$. One elementary proof is that
$P(T)-X$ is irreducible in $k(X)[T]$: its ideal in $k[T,X]$ has quotient
$k[T]$, and Gauss's lemma applies. If
$L=k(P(T),P(T^d))$, the just-proved minimum degree says
\[
 [L:k(P(T))]=M.
\]
The tower law then forces $L=k(T)$. Finally
$[k(T):k(P(T^d))]=dM$, which is the degree of the irreducible relation
in $X$; its degree in $Y$ is $M$.
\end{proof}

\begin{corollary}[An ordinary polynomial field identity]
\label{cor:polynomial-field}
Let $P\in k[T]\setminus k$, and let
$g=\gcd\{j\ge1:[T^j]P\ne0\}$. In characteristic zero,
\begin{equation}\label{eq:polynomial-field}
 k(P(T),P(T^d))=k(T^g)\qquad(d\ge2).
\end{equation}
\end{corollary}
\begin{proof}
Write $P(T)=Q(T^g)$, with $Q$ primitive in its exponents, and apply
\cref{thm:degree} to $Q$ in an ordinary Laurent Hahn field.
\end{proof}

This corollary can be read entirely in classical rational-function
algebra. We give it as a consequence and consistency check of the support
argument, not as a certified priority claim in polynomial decomposition
theory. The restriction to polynomials is real: for
$P(T)=T+T^{-1}$ one has $P(T^2)=P(T)^2-2$, and the generated field is
$k(T+T^{-1})$, a proper subfield of $k(T)$.

\begin{example}[A genuinely quadratic relation]
For $y=T^2+T$ and $d=2$, the primitive degree is two. The relation
\begin{align}\label{eq:sample-resultant}
 0={}&Y^2-(2X^2+6X+2)Y\\
     &+X^4+2X^3+2X^2\nonumber
\end{align}
is obtained by elimination and has the least possible $Y$-degree.
There is no rational function $R\in k(X)$ with $S_2y=R(y)$.
Thus the implicit algebraic classification is strictly broader than the
explicit rational-map classification in \cref{sec:rational}.
\end{example}

\section{Consequences for surreal and Gaussian omnific integers}
\label{sec:omnific}

\begin{theorem}[Omnific dilation classification]\label{thm:omnific}
Let $d\ge2$ and $x\in\Oz\setminus\Z$. The following are equivalent:
\begin{enumerate}[label=\textup{(\roman*)},leftmargin=2em,itemsep=2pt]
\item $x,S_dx$ are algebraically dependent over $\R$;
\item $x,S_dx$ are algebraically dependent over $\C$ inside $\SC$;
\item for some $\delta\in\No_{>0}$,
      $x=P(\omega^\delta)$ with $P\in\R[T]$ and $P(0)\in\Z$;
\item $x$ has finite normal-form support, and its positive exponents are
      commensurable.
\end{enumerate}
With the primitive exponent step, its least relation degree is $\deg P$.
For $x\in\Goz\setminus\Z[i]$, the identical classification holds with
$P\in\C[T]$, $P(0)\in\Z[i]$, and dependence over $\C$.
\end{theorem}
\begin{proof}
Choose a set-sized divisible subgroup of $\No$ containing the exponent
support, and use \eqref{eq:omega-bridge}. Apply
\cref{thm:negative,thm:degree}. In the real case, a complex polynomial
relation splits into its real and imaginary coefficient polynomials;
at least one is a nonzero real relation. Conversely a real relation is
a complex one. The coefficient at zero imposes exactly the integer or
Gaussian-integer condition and no condition on the other coefficients.
\end{proof}

\begin{corollary}[All distinct forward dilates]\label{cor:pairs}
If $x$ is an infinite-support omnific or Gaussian omnific integer, then
for every $d\ge2$ and every pair of integers $0\le a<b$,
\[
 S_d^a x\quad\text{and}\quad S_d^b x
\]
are algebraically independent over $\C$ (and over $\R$ in the real case).
\end{corollary}
\begin{proof}
The first member still has infinite support and the second is its
$S_{d^{b-a}}$ dilate. Apply \cref{thm:omnific}.
\end{proof}

\begin{example}[Finite but incommensurable]
The number
\[
 x=\omega^{\sqrt2}+\omega+1
\]
is an omnific integer, and $x,S_dx$ are algebraically independent for
every integer $d\ge2$: the positive exponents $1,\sqrt2$ are
$\Q$-linearly independent. This particular $x$ is also prime by the
published theorem of L'Innocente--Mantova \cite[Theorem~B]{LM}.
Its primality is cited, not proved or claimed as a contribution here.
The dilation independence is a different property.
\end{example}

\begin{example}[Comparable scales do not suffice]
Both
$\omega^{\sqrt2}+\omega$ and $\omega^{3/2}+\omega$
have two finite real exponents in the same Archimedean class of the
exponent group. Only the second is a polynomial in one monomial:
with $T=\omega^{1/2}$ it is $T^3+T^2$. Its least relation degree is
three; the first has no relation at all.
\end{example}

\begin{example}[One rational span, unbounded denominators]
The valid normal form
\[
 u_d=\sum_{n\ge0}\omega^{d^{-n}}
\]
has infinite support of order type $\omega$. All its exponents lie in
$\Q$, but not in any cyclic subgroup. It is algebraically independent
of $S_du_d$. In \cref{sec:order} we show that its first three transforms
satisfy an exact polynomial relation. Pairwise independence must not be
confused with independence of a longer tuple.
\end{example}

\begin{example}[A coefficient-field obstruction]
Let $a\in\R$ be transcendental over $\Q$ and $\delta>0$. Although
$x=a\omega^\delta$ satisfies
$S_dx=a^{1-d}x^d$ over $\R$, it has no first-order autonomous polynomial
relation over $\Q$. This follows from \cref{cor:coefficients}.
The constants allowed in a relation materially change the answer.
\end{example}

\section{Complete solutions for rational maps}\label{sec:rational}

We now classify the explicit equations $S_qy=R(y)$, with
$q\in\Q_{>1}$, $\Gamma$ divisible, and $R\in k(Z)$ nonconstant.
The polynomial relation theorem allows implicit ramification; in this
explicit case there is no such support ramification.

\begin{theorem}[Rational-map classification]\label{thm:rational}
Assume $k$ algebraically closed of characteristic zero and $\Gamma\ne0$
divisible. A nonconstant solution of
\begin{equation}\label{eq:rational}
 S_qy=R(y)
\end{equation}
exists exactly when $R$ has a fixed point $c\in\mathbb P^1(k)$ whose
local degree is $q$. In particular, $q$ must be an integer $d\ge2$.
For every such fixed point, let
\[
 M_c(Y)=Y-c\quad(c\in k),\qquad M_\infty(Y)=1/Y,
\]
and write
\[
 h=M_c\circ R\circ M_c^{-1}
   =A Z^d+O(Z^{d+1}).
\]
All nonconstant solutions associated with $c$ are exactly
\begin{equation}\label{eq:rational-solutions}
 y=M_c^{-1}\bigl(aH_a(t^\gamma)\bigr),
 \qquad \gamma\in\Gamma_{>0},\quad Aa^{d-1}=1,
\end{equation}
where $H_a$ is the normalized inverse B\"ottcher coordinate of
$a^{-1}h(aZ)$. Constant solutions are simply the finite fixed points of
$R$ where it is defined.
\end{theorem}
\begin{proof}
For $v(y)\ge0$, let $c$ be its residue. A pole of $R$ at $c$ would give
$v(R(y))<0$, impossible. Reducing the equation modulo the maximal ideal
therefore gives $R(c)=c$. For $v(y)<0$, the equation requires $\infty$
to be a fixed point. In either case $u=M_c(y)$ is nonzero of positive
valuation, and
\[
 S_qu=h(u).
\]
If $m\ge1$ is the local degree, comparison of valuations gives
$q\,v(u)=m\,v(u)$, hence $q=m$. Since $q>1$, this is an integer $d\ge2$.
The germ now has the form in \cref{prop:germ}, which gives all solutions
and their converse. Because $k$ is algebraically closed, the required
leading coefficients $a$ exist.
\end{proof}

In particular, not every nonlinear rational map admits a nonconstant
solution. The relevant property is a \emph{superattracting fixed point
of exactly the dilation degree}, not just a nonlinear term somewhere
in the formula. Over a non-algebraically-closed field the same description
applies whenever the fixed point and the required leading-coefficient
root lie in that field. For real solutions, these are genuine real-root
conditions, not automatic consequences of the complex theorem.

\begin{theorem}[Rational-map rigidity in the negative-support ring]
\label{thm:rational-negative}
Let $\operatorname{char}k=0$, $y\in\A\setminus k$, and $d\ge2$.
If $S_dy=R(y)$ for $R\in k(Z)$, then
\begin{equation}\label{eq:affine-monomial}
 y=\alpha t^{-\gamma}+b,
 \qquad
 R(Z)=\alpha^{1-d}(Z-b)^d+b,
\end{equation}
with $\alpha\in k^\times$, $b\in k$, and $\gamma\in\Gamma_{>0}$.
Conversely every such pair satisfies the equation.
For omnific solutions $b\in\Z$, and for Gaussian omnific solutions
$b\in\Z[i]$.
\end{theorem}
\begin{proof}
Writing $R=A/B$ in reduced form, the equation
$B(y)S_dy-A(y)=0$ has degree one in the dilated variable. The denominator
is nonzero since $y$ is transcendental over $k$. By the degree bound in
\cref{thm:negative}, $y$ is affine in one monomial, giving the first
formula. Substitute this into the equation, and use transcendence of
the monomial to obtain the displayed rational-function identity.
The converse is direct substitution. The constant-term restrictions
are exactly the definitions of the two omnific rings.
\end{proof}

\begin{corollary}[Simultaneous rational equations]\label{cor:simultaneous}
Suppose the same nonconstant $y\in\A$ satisfies
$S_py=P(y)$ and $S_qy=Q(y)$ for integers $p,q\ge2$ and rational functions
$P,Q$. Then, with one common $\alpha,b,\gamma$,
\[
 y=\alpha t^{-\gamma}+b,\quad
 P(Z)=\alpha^{1-p}(Z-b)^p+b,\quad
 Q(Z)=\alpha^{1-q}(Z-b)^q+b.
\]
In particular, $P$ and $Q$ commute. Even without the negative-support
restriction, simultaneous nonconstant solutions force $P\circ Q=Q\circ P$.
\end{corollary}
\begin{proof}
Apply \cref{thm:rational-negative}; the representation is unique from
its leading monomial and constant term. For the last assertion, commute
the dilations and use the fact that they fix the coefficients:
$(P\circ Q)(y)=(Q\circ P)(y)$. Transcendence of $y$ turns equality at
$y$ into equality of rational functions.
\end{proof}

\section{A sharp infinitesimal-tail obstruction}\label{sec:tail}

The rational-map obstruction has a quantitative form. It identifies the
first infinitesimal term that prevents an otherwise legitimate surreal
solution from being an omnific integer.

\begin{theorem}[First forbidden tail]\label{thm:tail}
Let $P\in\C[Z]$ have degree $d\ge2$, with leading coefficients
$p_d,p_{d-1}$. Put
\[
 b=-\frac{p_{d-1}}{d p_d},\qquad
 p_d\alpha^{d-1}=1,
 \qquad Q(X)=\alpha^{-1}(P(\alpha X+b)-b).
\]
Then $Q$ is monic and centered. If $Q\ne X^d$, write
\[
 Q(X)=X^d+cX^{d-r}+\text{lower powers},
 \qquad 2\le r\le d,\quad c\ne0.
\]
For each positive scale $\gamma$, the infinite solution with leading
coefficient $\alpha$ has expansion, with $T=t^{-\gamma}$,
\begin{equation}\label{eq:tail}
 y=\alpha T+b-\frac{\alpha c}{d}T^{1-r}+O(T^{-r}).
\end{equation}
Its first positive-valuation term has valuation exactly $(r-1)\gamma$.
In particular it does not belong to the negative-support ring.
\end{theorem}
\begin{proof}
The normalization eliminates the coefficient of $X^{d-1}$ and makes the
leading coefficient one. The infinity case of \cref{thm:rational}
gives a unique normalized Laurent solution
\[
 F(T)=T+\sum_{n\ge0}b_nT^{-n},\qquad F(T^d)=Q(F(T)),
\]
with $y=\alpha F(T)+b$. Comparison at degree $d-1$ gives $b_0=0$.
Successive comparison at degrees $d-2,d-3,\ldots,d-r+1$ gives
$b_1=\cdots=b_{r-2}=0$. At degree $d-r\ge0$, the left side has no
term, and the right-side coefficient is $d b_{r-1}+c$. Therefore
$b_{r-1}=-c/d\ne0$. This proves the expansion and the exact valuation.
\end{proof}

Here $O(T^{-r})$ denotes an ordinary formal Laurent tail with powers at
most $-r$; it is not an analytic estimate at all surreal magnitudes.
After $T=t^{-\gamma}$, this notation has the exact support meaning
required by the theorem.

\subsection{Examples and exact coefficients}

For $P(Z)=Z^2+1$, the unique infinite branch at each positive scale begins
\begin{equation}\label{eq:quadratic-example}
 y=T-\frac12T^{-1}-\frac38T^{-3}-\frac3{16}T^{-5}
      -\frac{45}{128}T^{-7}+\cdots.
\end{equation}
It is a perfectly well-defined real surreal solution of
$S_2y=y^2+1$ after $T=\omega^\gamma$, but is not omnific.
There are no nonconstant omnific solutions to this equation.

For $P(Z)=Z^2-2$ there is an exact finite Laurent identity:
\[
 y=T+T^{-1},\qquad S_2y=T^2+T^{-2}=y^2-2.
\]
Again the infinitesimal term prevents an omnific solution. This example
also explains why the full Hahn-field conclusion is Laurent-series
rigidity, not polynomial rigidity.

For $P(Z)=Z^3+Z$ one obtains
\[
 y=T-\frac13T^{-1}-\frac8{81}T^{-5}-\frac8{243}T^{-7}+\cdots.
\]
By contrast, $P(Z)=(Z-3)^2+3$ has the omnific solutions
$y=3+\omega^\gamma$, since its centered normalization is exactly $Z^2$.
The finite exact checks supplied with the article verify the displayed
coefficients and the coordinate identities through the stated truncation
orders; \cref{thm:tail} proves the infinite statement.

\section{A strict hierarchy of autonomous algebraic orders}
\label{sec:order}

\begin{definition}
For a fixed dilation $S_d$, define the \emph{autonomous algebraic order}
of $y$ over $k$ to be the least $r\ge0$ for which
\[
 y,S_dy,\ldots,S_d^r y
\]
are algebraically dependent over $k$. If no such $r$ exists, set the
order to $+\infty$. Constants have order zero. This definition concerns
algebraic dependence of a transform tuple; it is not the usual order of
an arbitrary variable-coefficient linear equation.
\end{definition}

\begin{theorem}[An infinite-support element of exact order two]
\label{thm:order-two}
For each integer $d\ge2$, the omnific integer
\[
 u_d=\sum_{n\ge0}\omega^{d^{-n}}
\]
has exact autonomous algebraic order two over $\C$ for $S_d$. It satisfies
\begin{equation}\label{eq:order-two}
 S_d^2u_d-S_du_d=(S_du_d-u_d)^d.
\end{equation}
Every two distinct forward transforms are nevertheless algebraically
independent over $\C$.
\end{theorem}
\begin{proof}
Its decreasing normal-form exponents make the sum legitimate. Dilation
is strong, and direct comparison of supports gives
\[
 S_du_d=u_d+\omega^d.
\]
Applying $S_d$ once more gives
$S_d^2u_d-S_du_d=\omega^{d^2}=(\omega^d)^d$, proving
\eqref{eq:order-two}. The polynomial
$Z-Y-(Y-X)^d$ is nonzero, so the order is at most two.
By \cref{thm:omnific}, the infinite support makes $u_d,S_du_d$
algebraically independent. Thus the order is exactly two.
The last assertion is \cref{cor:pairs}.
\end{proof}

Decreasing-denominator generalized series and their linear Mahler
identities are already present in the literature; compare
\cite[Remark~2.18]{CDDM} and \cite{GG}. The series pattern is therefore
not offered as a new construction. The point here is the rigorous
\emph{minimal autonomous order} and the resulting sharp boundary of the
first-order classification in the omnific setting.

\begin{theorem}[Every finite order occurs]\label{thm:all-orders}
Let $g_1,\ldots,g_r$ be positive surreal exponents linearly independent
over $\Q$, and let $a_1,\ldots,a_r\in\R^\times$. Then
\[
 x=\sum_{j=1}^r a_j\omega^{g_j}
\]
has exact autonomous algebraic order $r$ over $\C$ for every integer
$d\ge2$.
\end{theorem}
\begin{proof}
Put $T_j=\omega^{g_j}$. The $T_j$ are algebraically independent over
$\C$, since distinct monomials in them have distinct normal-form
exponents. For $0\le i<r$, put
\[
 Y_i=\sum_{j=1}^r a_jT_j^{d^i}.
\]
Their Jacobian determinant is
\begin{equation}\label{eq:jacobian}
 \left(\prod_{j=1}^r a_j\right)
 \left(\prod_{i=0}^{r-1}d^i\right)
 \det\bigl(T_j^{d^i-1}\bigr)_{0\le i<r,\,1\le j\le r}.
\end{equation}
The last determinant is nonzero: its permutation terms have distinct
exponent vectors, since the integers $d^i-1$ are distinct.
In characteristic zero, a nonzero Jacobian implies algebraic
independence. For completeness, a polynomial relation of least total
degree would, on differentiation, yield a nontrivial linear relation
between the differentials $dY_i$. At least one evaluated partial
derivative is nonzero, because otherwise a nonzero partial derivative
would be a relation of smaller degree. This contradicts the determinant.

Thus the first $r$ transforms are independent. The first $r+1$ lie in
the field $\C(T_1,\ldots,T_r)$ of transcendence degree $r$, so they are
dependent. The least order is exactly $r$.
\end{proof}

\begin{proposition}[Bounds from the rational rank of finite support]
\label{prop:rank-bounds}
Let $x$ be a nonconstant finite-support omnific integer and let $r$ be
the dimension over $\Q$ of the span of its positive exponents. Its
order is at most $r$. It is exactly one for $r=1$ and exactly two for
$r=2$; for $r\ge3$, the present results give the bounds $2\le\text{order}\le r$.
\end{proposition}
\begin{proof}
Choose a rational basis of the exponent span and clear the finitely many
denominators. Then every transform of $x$ lies in a rational function
field generated by $r$ algebraically independent monomials. This gives
the upper bound. The first-order classification gives precisely the
remaining assertions.
\end{proof}

\section{Boundaries that the conclusions cannot cross}\label{sec:boundaries}

\subsection{Positive characteristic}
If $k=\mathbb F_p$ and the dilation is $S_p$, then
\[
 S_py=y^p
\]
for every Hahn series $y$: Frobenius acts coefficientwise on these
strong sums. Hence completely arbitrary well-ordered supports can satisfy
a degree-one autonomous relation in the dilated variable. The
characteristic-zero hypothesis cannot be dropped. In the proof, both
Newton--Puiseux and the nonvanishing of the coefficient $d$ are substantive.

\subsection{A named nonconstant coefficient}
The infinite-support series $u_d$ satisfies the linear equation
\[
 S_du_d-u_d=\omega^d.
\]
The right side is not in the ordinary constant field. Thus there is no
extension of the support-collapse theorem to unrestricted nonconstant
coefficients. The distinction from variable-coefficient Mahler algorithms
\cite{FR,CDDM} is structural, not a matter of proof technique.

\subsection{Identity dilation and arbitrary rational dilations}
For $d=1$, every series satisfies $S_1y-y=0$, so the theorem is false.
For a nonintegral rational $q>1$, \cref{thm:rational} rules out explicit
rational-map solutions. That argument does not classify implicit
relations $F(y,S_qy)=0$. In the Puiseux step the leading order becomes
$qe$, and taking an $e$th root need not produce an ordinary integer-order
germ. The proof given here must not be quoted as covering that case.

\subsection{Full first-order theories}
The decision procedure below is for one specified autonomous equation
with algebraic-number coefficients. It says nothing about arbitrary
quantifier alternations in an expanded surreal field, arbitrary systems
of omnific Diophantine equations, or equality of arbitrary infinite
normal-form descriptions. In particular, it does not contradict the
undecidability phenomena discussed in the repository's single-dilation
report \cite{RepoDilation}.

\section{Effective finite shapes and a decision procedure}
\label{sec:effective}

The sharp degree bound turns the omnific problem into a finite algebraic
problem, even though a general omnific integer is specified by a
potentially transfinite normal form.

\begin{theorem}[Finitely many bounded polynomial shapes]
\label{thm:finite-polynomials}
Let $k$ be algebraically closed of characteristic zero,
$0\ne F\in k[X,Y]$, and $d\ge2$. Put $D=\deg_YF$.
The set
\begin{equation}\label{eq:polynomial-shapes}
 \mathcal P_{\le D}(F,d)=
 \{P\in k[T]:1\le\deg P\le D,
          \ F(P(T),P(T^d))=0\}
\end{equation}
is finite. Every nonconstant solution in a negative-support Hahn ring
has the form $P(t^{-\delta})$ with $P$ in this finite set and $\delta>0$.
The analogous statement holds elementwise for omnific and Gaussian
omnific integers, retaining the appropriate coefficient and constant-term
conditions.
\end{theorem}
\begin{proof}
Fix a degree $M$. For a polynomial $P$ of that degree satisfying the
identity, view $P(t^{-1})$ in $k((t^\Q))$. It has valuation $-M$.
By \cref{thm:profiles} it equals $f(t^\delta)$ for one of finitely many
infinity profiles $f$ of pole order $e$. Valuation comparison forces
$\delta=M/e$. Thus, for each fixed $M$ and profile, there is at most
one possible polynomial $P$. Finiteness follows on taking
$1\le M\le D$. Exhaustion of all negative-support solutions is precisely
the degree-bounded assertion of \cref{thm:negative} and transcendence of
the chosen monomial. Filtering the coefficients gives the omnific versions.
\end{proof}

The bound on the polynomial degree is essential to this finiteness
statement. For example, $S_dy=y^d$ is satisfied by $P(T)=T^n$ for
arbitrarily large $n$. These do not create new shapes modulo the free
monomial scale, since $(t^{-\delta})^n=t^{-n\delta}$.

\begin{theorem}[Decidable nonconstant omnific solvability]
\label{thm:decidable}
Let $F\in\overline{\Q}[X,Y]\setminus\{0\}$ be supplied by exact
algebraic-number coefficients and let $d\ge2$. There is an algorithm to
decide whether $F(x,S_dx)=0$ has a nonconstant solution in $\Oz$, or in
$\Goz$. It produces a finite list of polynomial shapes whose positive
surreal monomial evaluations give all such solutions. Every coefficient
of those shapes is an algebraic number.
\end{theorem}
\begin{proof}
If $D=\deg_YF=0$, no nonconstant solution exists. Otherwise, for each
$1\le M\le D$, introduce unknowns $a_0,\ldots,a_M$ and form
\[
 P(T)=a_0+a_1T+\cdots+a_MT^M.
\]
Set every coefficient in $T$ of $F(P(T),P(T^d))$ equal to zero and impose
$a_M\ne0$, for instance by adding an auxiliary inverse variable.
These are finite polynomial systems with algebraic-number coefficients.
By \cref{thm:finite-polynomials}, each has only finitely many solutions
in $\C$. A finite algebraic set over $\overline\Q$ has algebraic-number
coordinates. Its points can be computed by standard zero-dimensional
polynomial-system methods. Alternatively, separate real and imaginary
parts and use exact real-algebraic point isolation \cite{BPR}; a nonzero
complex leading coefficient is imposed by inverting its squared modulus.

For $\Oz$, retain precisely the shapes with every coefficient real and
$a_0\in\Z$. For $\Goz$, retain those with $a_0\in\Z[i]$.
Reality, equality, and integrality of exact algebraic numbers are
decidable. The retained list is nonempty exactly when a nonconstant
solution exists. Its evaluations at $T=\omega^\delta$ give all solutions
by \cref{thm:omnific}. A witness may always be produced already at
$T=\omega$.
\end{proof}

The algorithm is a termination result, not a practical complexity claim.
The coefficient systems can be large, and a production implementation
would exploit Puiseux branches and the finite-profile structure before
using general elimination. The supplied Python program implements exact
checks of representative identities, not this complete elimination
algorithm.

\subsection{A finite computational specification}

The equation expanded in the proof has degree at most
$M(\deg_XF+d\deg_YF)$ in $T$. Thus there is an explicit finite bound on
the number of coefficient equations. Every returned candidate is
certified simply by the exact polynomial identity
$F(P(T),P(T^d))=0$, the leading-coefficient test, and the constant-term
condition. No truncation of an unknown infinite surreal number is used
as a certificate of support termination.

The Newton-weight test of \cref{cor:existence} is a different algorithm:
it decides nonconstant solvability in the full surcomplex field, where
infinite Laurent tails are allowed. The two procedures can therefore
legitimately give different answers, as for $S_2x=x^2+1$.

\section{Further research questions and priorities}\label{sec:questions}

The questions below are proposed research directions arising from this
article. Unless explicitly stated otherwise, they are not claimed to be
previously published named open problems. Their purpose is to specify
concrete next theorems rather than to suggest that the present first-order
classification covers a wider theory than it actually does.

\begin{question}[Exact order from finite support; high priority]
For a finite-support omnific integer whose positive exponents have rational
rank $r\ge3$, is its autonomous order always $r$? The answer is established
here for $r\le2$ and for sums supported on a rationally independent set.
The unresolved case allows several supported monomials in the same
$r$-dimensional exponent span. A useful first target is rank three with
four support terms. The bound in \cref{prop:rank-bounds} isolates the
precise remaining gap.
\end{question}

\begin{question}[The support spectrum at order two; high priority]
Which well-ordered Hahn supports can occur for omnific solutions of
$F(x,S_dx,S_d^2x)=0$ over the ordinary constants? Infinite support of
order type $\omega$ occurs by \cref{thm:order-two}. Determine whether
fixed equation degree or fixed autonomous order bounds the possible
support order types, denominator growth, or rational rank. The goal is
a classification or a sharp no-bound theorem, not an extrapolation from
the one telescoping example.
\end{question}

\begin{question}[Formal verification of the core theorem; high priority]
Formalize the chain from strong Hahn evaluation through the monomial
recognizer, the formal B\"ottcher coordinate, and the evaluated Puiseux
branch to \cref{thm:collapse}. The most substantial imported obligation
is a usable Newton--Puiseux theorem with the \emph{individual branch}
degree bound. A staged formalization can state that theorem as an explicit
parameter before replacing it with a verified construction; it should
not silently add it as an axiom and call the entire result verified.
\end{question}

\begin{question}[An efficient shape solver]
Implement \cref{thm:decidable} for rational or algebraic coefficients,
using the Newton-weight filter, Puiseux branches, and polynomial degree
bounds before general elimination. Prove bounds in terms of the input
bidegree, coefficient height, and $d$. Can the finite candidate list be
computed without constructing any nonpolynomial B\"ottcher profile?
Can equivalent shapes under monomial rescaling be canonically minimized?
\end{question}

\begin{question}[Nonintegral rational dilation]
Classify autonomous implicit equations $F(y,S_qy)=0$ for
$q\in\Q_{>1}\setminus\Z$. The explicit rational case has no nonconstant
solutions, but a ramified correspondence can have the required rational
Newton slope. Determine whether an appropriate generalized B\"ottcher
coordinate still forces a controlled support subgroup, or permits
unbounded denominators already at order one.
\end{question}

\begin{question}[Tame positive-characteristic variants]
When $\operatorname{char}k=p>0$ and $p\nmid d$, which support-collapse
statements survive under restrictions on Puiseux ramification? Separate
failure caused by wild algebraic branches from failure of the coordinate
construction. The Frobenius example with $d=p$ shows that no unrestricted
analogue is possible, but does not settle the tame cases.
\end{question}

\begin{question}[Systems and higher-dimensional correspondences]
For a tuple $\mathbf y$ satisfying
$\mathbf S_d\mathbf y=\mathbf R(\mathbf y)$, what bounds the number of
independent support scales? In one dimension, a superattracting germ is
formally conjugate to a monomial. Higher-dimensional germs need not have
such a coordinate. Identify hypotheses under which a finite-rank support
module, rather than one cyclic group, is forced.
\end{question}

\begin{question}[Controlled nonconstant coefficients]
Replace the ordinary constant field by a finitely generated monomial field
$k(t^{\Delta})$. A linear equation already admits unbounded denominators,
so cyclic support is too strong. Find the correct substitute: perhaps a
bound on the support module generated under multiplication and division
by $d$, relative to $\Delta$, together with an ordinal support bound.
Distinguish finite-generation assertions from mere countability.
\end{question}

\begin{question}[Arithmetic and convergence of the profiles]
For algebraic coefficients, all profile coefficients are algebraic by
\cref{cor:coefficients}. Determine effective denominator and height bounds,
and compare their Archimedean and non-Archimedean analytic convergence.
The underlying B\"ottcher arithmetic is already an established subject
\cite{SS}; the new target is a uniform statement through the ramified
correspondence and the finite-profile list.
\end{question}

\begin{question}[Simultaneous algebraic correspondences]
Describe the intersection of the finite-profile solution classes for two
multiplicatively independent dilations. The rational omnific case is
completely rigid by \cref{cor:simultaneous}, but implicit correspondences
and full Laurent tails allow more geometry. Determine which compatibility
conditions replace commutation of two rational maps.
\end{question}

\begin{question}[Definability of bounded shape complexity]
For each fixed bound $M$, describe the locus of omnific elements whose
least first-order relation degree is at most $M$, in the language with
one named dilation and in weaker languages. The repository already
recovers normal-form data from a dilation \cite{RepoDilation}; the
question here is the complexity and uniformity of the \emph{bounded-degree}
stratification. Do not confuse a countable union of definable loci with
one first-order definable set.
\end{question}

\begin{question}[Interaction with surreal exponential and derivation]
The dilation in this article is defined by the normal-form exponent group,
not by the surreal exponential. Study how the finite-profile classes
behave under the surreal exponential and under specified surreal
exponential-compatible derivations. Even a negative closure theorem would
clarify which parts of difference algebra and differential algebra can be
combined without changing the constant field or losing support control.
\end{question}

The recommended order is to formalize the support-collapse theorem,
implement the finite-shape solver, and then attack the rank-three and
order-two problems. Each has a precise success criterion and can be
pursued without a general theory of analysis on the entire proper class.

\section{Proof audit and concluding assessment}\label{sec:audit}

\subsection{Dependency ledger}

\begin{center}
\small
\begin{tabularx}{\textwidth}{@{}p{0.30\textwidth}Y@{}}
\toprule
Input or conclusion & Status and dependency \\
\midrule
Hahn field arithmetic and normal forms & Classical framework; used in
set-sized workspaces. \\
Strong support lemma & Classical Neumann lemma; supplies all infinite
formal evaluations. \\
Newton--Puiseux factorization & Classical characteristic-zero theorem;
individual denominator bound proved here by the roots-of-unity orbit. \\
Formal B\"ottcher coordinate & Classical theorem; explicit logarithmic
construction and uniqueness proof included. \\
Monomial recognizer & Already present for $S_2$ in the repository;
general-$d$ least-term proof included. \\
Support collapse and finite profiles & Proposed contribution; deductions
from the preceding inputs, with no rank-one hypothesis. \\
Omnific criterion and sharp degree & Proposed consequences; follow by
intersecting cyclic support with a bounded interval. \\
Finite-shape decision procedure & Uses the sharp degree bound, finite
profiles, and standard exact finite algebraic-system solving. \\
Higher-order examples & Exact written identities and algebraic
independence arguments; the generalized series pattern itself is prior. \\
\bottomrule
\end{tabularx}
\end{center}

The most delicate proof points have been made explicit: the selected
Puiseux denominator is not a common splitting denominator; the root
of the algebraic branch must be chosen to equal the actual $S_dz$;
formal evaluation must respect composition and dilation; and the
nondivisible-group descent uses a subgroup intersection containing the
leading valuation. None of these is delegated to numerical evidence.

\subsection{What has and has not been checked}

The accompanying exact checks cover six inverse B\"ottcher expansions
through degree 32 in the normalized power-series variable, twenty
resultant examples, ten nonzero Jacobian instances, forty-five finite
telescoping identities, four Frobenius examples, and six Newton-weight
examples. All passed in the execution recorded in
\texttt{verification\_results.json}. The calculations use exact rational
or polynomial arithmetic. Their purpose is to detect sign, normalization,
indexing, and small-example errors.

These checks do not certify the full support theorem, the correctness
of every classical dependency, or the decision procedure's implementation.
No Lean build was performed and no new Lean theorem is supplied. The
repository's pre-existing formalization status is not inherited by this
article. Independent mathematical review is still needed.

The source search covered the named primary literature and the specified
repository reports and inventory. The Nishioka--Nishioka paper was located
and its abstract-level scope checked; its full text was not available
through the retrieval route used here, and no detailed theorem from it
is invoked. This is a limitation on novelty assessment, not a missing
premise in any proof. An absence of matching search results is not proof
of priority.

\subsection{Conclusion}

The main structural distinction is now exact. At autonomous order one,
an algebraic relation with an integer exponent dilation collapses an
arbitrary Hahn support to a cyclic Laurent scale. In the omnific ring,
that collapse forces a finite polynomial in one monomial, and the
primitive polynomial degree is recoverable as an algebraic extension
degree. For a fixed equation there are only finitely many shapes modulo
scale, giving an effective solvability reduction for algebraic input.

At order two the finite-support conclusion already fails. This sharp
change is a useful organizing principle for further surreal difference
algebra. It supplies both positive classifications and explicit negative
results without relying on a problematic global topology or treating a
proper class as an ordinary analytic field.

\appendix
\section{A worked elimination identity}\label{app:elimination}

Let $X=T^2+T$ and $Y=T^4+T^2$. From $T^2=X-T$ one obtains
\[
 T^4=(X-T)^2=X^2+X-(2X+1)T,
\]
and hence
\[
 Y=X^2+2X-2(X+1)T.
\]
Therefore
\[
 T=\frac{X^2+2X-Y}{2(X+1)}.
\]
The denominator is nonzero in $k(T)$ in characteristic zero. Substitution
into $T^2+T-X=0$ and clearing denominators gives
\[
 Y^2-(2X^2+6X+2)Y+X^4+2X^3+2X^2=0.
\]
This simultaneously exhibits the relation and reconstructs the primitive
monomial from the pair. The degree theorem proves minimality; the
calculation alone would not rule out a hidden linear relation.

\section{Reproducing the finite checks}\label{app:checks}

The file \texttt{verify.py} requires Python 3.9 or later and SymPy.
It uses only exact arithmetic. Run
\begin{verbatim}
python -m pip install sympy
python verify.py
\end{verbatim}
The second command writes \texttt{verification\_results.json} beside the
script and prints the same report. The recorded run used Python 3.13.5
and SymPy 1.14.0. The script is not an implementation of infinite Hahn
series and does not pretend that truncation is an infinite proof.

For the inverse B\"ottcher checks, set $L(z)=zF(1/z)$ for a monic centered
polynomial $Q(X)=\sum q_jX^j$ of degree $d$. The checked identity is
\[
 L(z^d)=\sum_{j=0}^d q_j z^{d-j}L(z)^j.
\]
The coefficient of degree $m$ on the right contains $d[ z^m ]L$ plus
already known coefficients. Since $d\ne0$, this is an exact recursion.
The program independently recomputes the two sides after constructing
the coefficients and verifies equality through degree 32.

For the finite sums $u_{d,N}=\sum_{n=0}^N\omega^{d^{-n}}$, the checked
identity is deliberately the correct truncated one,
\[
 S_du_{d,N}-u_{d,N}=\omega^d-\omega^{d^{-N}},
\]
not the infinite identity with its boundary term silently deleted.
The infinite identity follows separately from coefficientwise comparison
of the full well-ordered supports.

\begin{thebibliography}{99}
\small\raggedright
\bibitem{Repo}
Vladimir Reshetnikov and contributors,
\emph{Surreal}, repository and research inventory, inspected at commit
\texttt{b895e8672990e8b5a56f97dd7246f9dd5f86c771}, 23 September 2026.
\url{https://github.com/VladimirReshetnikov/Surreal}.
Relevant inventory: \texttt{docs/README.md} and \texttt{docs/manifest.tex}.

\bibitem{RepoDilation}
\emph{A Single Dilation Recovers Hahn Support: Exact difference cohomology,
definable normal forms, and centralizers of monomial automorphisms}, research report in \repo, 22 September 2026.
\texttt{docs/surcomplex/single-dilation-hahn-support/article.tex},
at the commit specified in \cite{Repo}. AI-assisted, unrefereed repository
manuscript; cited as prior repository work, not as peer-reviewed literature.

\bibitem{RepoDynamics}
\emph{Surcomplex Dynamics: Linearization Thresholds, Periods, and Normal
Forms}, merged research report in \repo, 21--22 September 2026.
\texttt{docs/surcomplex/dynamics-and-normal-forms/article.tex},
at the commit specified in \cite{Repo}. AI-assisted, unrefereed manuscript.

\bibitem{Neumann}
B.~H.~Neumann, On ordered division rings,
\emph{Transactions of the American Mathematical Society} \textbf{66}
(1949), 202--252.
\url{https://doi.org/10.1090/S0002-9947-1949-0032593-5}.

\bibitem{MC}
Bassel Mannaa and Thierry Coquand,
\emph{Dynamic Newton--Puiseux Theorem}, arXiv:1304.6770v2 (2013).
\url{https://arxiv.org/abs/1304.6770}.

\bibitem{SS}
Adriana Salerno and Joseph H.~Silverman,
\emph{Integrality properties of B\"ottcher coordinates for one-dimensional
superattracting germs}, arXiv:1708.09275v2 (2017), especially Proposition~1.
\url{https://arxiv.org/abs/1708.09275}.

\bibitem{NN}
Kumiko Nishioka and Seiji Nishioka,
Autonomous equations of Mahler type and transcendence,
\emph{Tsukuba Journal of Mathematics} \textbf{39}(2), 251--257;
issue year 2015, online publication 2016.
\url{https://doi.org/10.21099/tkbjm/1461270059}.
Scope checked from the abstract; no detailed result is used as a premise.

\bibitem{GG}
R.~Gontsov and I.~Goryuchkina,
On the existence and convergence of formal power series solutions of
nonlinear Mahler equations,
\emph{Journal of Symbolic Computation} \textbf{128} (2025), 102399.
\url{https://doi.org/10.1016/j.jsc.2024.102399}.

\bibitem{CDDM}
Fr\'ed\'eric Chyzak, Thomas Dreyfus, Philippe Dumas, and Marc Mezzarobba,
\emph{Computing solutions of linear Mahler equations},
arXiv:1612.05518v2 (2018), especially Remark~2.18.
\url{https://arxiv.org/abs/1612.05518}.

\bibitem{FR}
Colin Faverjon and Julien Roques,
\emph{Hahn series and Mahler equations: Algorithmic aspects},
arXiv:2412.04928 (2024).
\url{https://arxiv.org/abs/2412.04928}.

\bibitem{FP}
Colin Faverjon and Marina Poulet,
\emph{Computing basis of solutions of any Mahler equation},
arXiv:2511.18877 (2025).
\url{https://arxiv.org/abs/2511.18877}.
Abstract-level contextual citation only.

\bibitem{FRPurity}
Colin Faverjon and Julien Roques,
\emph{A purity theorem for Mahler equations},
arXiv:2512.10661 (2025).
\url{https://arxiv.org/abs/2512.10661}.
Abstract-level contextual citation only.

\bibitem{LM}
Sonia L'Innocente and Vincenzo Mantova,
A factorisation theory for generalised power series and omnific integers,
\emph{Advances in Mathematics} \textbf{442} (2024), 109513.
\url{https://arxiv.org/abs/1710.07304}; version 5, 22 January 2024.

\bibitem{BPR}
Saugata Basu, Richard Pollack, and Marie-Fran\c{c}oise Roy,
\emph{Algorithms in Real Algebraic Geometry}, second edition,
Springer, 2006.
\url{https://doi.org/10.1007/3-540-33099-2}.

\bibitem{Wiki}
Wikipedia contributors, \emph{Surreal number}, consulted 23 September 2026.
\url{https://en.wikipedia.org/wiki/Surreal_number}.
Orientation only; not an input to the proofs.
\end{thebibliography}
\end{document}
